\section{Attempts reliability and selection}\label{sec:C26:attempts}
Finding one successful answer among several attempts differs from succeeding reliably on repeated execution. Under independent identical success probability $p$, the former is $1-(1-p)^k$ and the latter is $p^k$. These formulas are illustrative models, not general laws for agent systems. Attempts can share task difficulty, memory, environment state, or sampling behavior, making their joint distribution substantially different.

For $n$ sampled candidates with $c$ successes, the familiar combinatorial pass-at-$k$ estimator is $1-\binom{n-c}{k}/\binom{n}{k}$ when the sampling and interpretation conditions apply~\cite{chen2021codeeval}. The domain is $1\leq k\leq n$ and $0\leq c\leq n$, all integers, with $\binom{n-c}{k}=0$ when $k>n-c$. It estimates discovery of a success among selected candidates, not a deployable selector's performance. An actual system needs a gate that identifies a valid candidate without access to hidden labels. Selector false acceptance and false rejection belong inside end-to-end evaluation.

For a fixed realized candidate set, that expression is exactly the fraction of uniformly chosen size-$k$ subsets containing a success. Its interpretation as an unbiased estimate of fresh $k$-attempt discovery additionally needs the corresponding exchangeable sampling model; adaptive revisions are not automatically such samples. Task heterogeneity introduces a separate distinction. If attempts are independent conditional on task $q$ with success probability $p(q)$, then
\begin{equation}
 \theta_{\mathrm{any},k}=\mathbb E_q[1-(1-p(q))^k],\qquad
 \theta_{\mathrm{all},k}=\mathbb E_q[p(q)^k].
\end{equation}
For integer $k\geq1$ and $\bar p=\mathbb E_q[p(q)]$, Jensen's inequality gives $\theta_{\mathrm{any},k}\leq1-(1-\bar p)^k$ and $\theta_{\mathrm{all},k}\geq\bar p^k$. Thus even conditional independence does not justify raising a pooled mean to a power. These are analytical comparisons under the stated model, not estimates of any cited benchmark's reliability.

Attempt budgets should count failed parsing, tool errors, and abandoned generations where they consume resources. The mini-swe-agent accounting change discussed in Chapter~\ref{ch:C23} illustrates why this is implementation-sensitive. Equal nominal attempts can hide unequal token use or paid failed calls. Report measured cost and consumed resources rather than only configured maxima.

Repeated-run reliability should retain the benchmark's exact definition. Some protocols ask whether every repetition succeeds; others measure average success or consistency of a required constraint. A paper using one definition cannot be compared directly with another merely because both use a superscript or similar label. State the aggregation rule in ordinary language and provide the underlying per-task outcomes when possible.

Adaptive selection increases the effective search budget. A system that repeatedly revises prompts, harnesses, or memory using benchmark feedback has used the benchmark as an optimization resource. RRSI explicitly addresses overfitting in harness evolution~\cite{rrsi2026}. The evaluation should preserve an untouched assessment set and record candidate counts, rejected changes, and evaluator access. Frozen weights do not imply absence of learning from test feedback.

An untouched confirmation set is one design; certified reuse requires a different protocol. REUSE permits feedback-adaptive proposals only through promotion decisions from an initially independent evaluation sample and allocates error across all possible promotion histories~\cite{ch02_reuse2026}. Returning detailed scores or failure traces to the proposer breaks that theorem's decision-only restriction unless the extra transcript is accounted for. A generic confidence sequence protects optional stopping under its assumptions; it does not erase information already leaked into candidate construction. The public report should therefore describe the feedback channel as precisely as the tested artifact.

